\begin{examplebox}
	\textbf{\textcolor{CExample}{例 1（基础）}}

	\textbf{\textcolor{CExample}{题目：}}
	已知抛物线$y^2=6x$上有一点$P\left(\dfrac{3}{2},3\right)$，求过点$P$的切线方程。

	\textbf{\textcolor{CExample}{解题步骤：}}
	\begin{description}[style=nextline, leftmargin=2.8em, labelsep=0.5em, itemsep=0.45em]
		\item[\textbf{第一步（验证点在线上）：}] 验证$3^2=9=6\times\dfrac{3}{2}$，故点$P$在抛物线上。
			\par\textbf{理由：}切线公式要求切点在抛物线上，使用前必须验证。
		\item[\textbf{第二步（确定参数）：}] 由$y^2=6x$得$2p=6$，故$p=3$。
			\par\textbf{理由：}公式$y_0 y = p(x+x_0)$中的$p$对应抛物线方程$y^2=2px$中的$p$。
		\item[\textbf{第三步（套用公式）：}] 将$y_0=3$，$x_0=\dfrac{3}{2}$，$p=3$代入$y_0 y = p(x+x_0)$得$3y = 3\left(x+\dfrac{3}{2}\right)$，即$y = x+\dfrac{3}{2}$。
			\par\textbf{理由：}直接二级结论，一步得到切线方程。
	\end{description}

	\textbf{\textcolor{CExample}{关键结论：}}
	所求切线方程为$y = x + \dfrac{3}{2}$（也可写作$2y=2x+3$）。

	\medskip
	\textbf{\textcolor{CExample}{例 2（稍变形）}}

	\textbf{\textcolor{CExample}{题目：}}
	过点$A(4,5)$作抛物线$y^2=4x$的切线，求切线方程。

	\textbf{\textcolor{CExample}{解题步骤：}}
	\begin{description}[style=nextline, leftmargin=2.8em, labelsep=0.5em, itemsep=0.45em]
		\item[\textbf{第一步（设切点坐标）：}] 设切点为$P(x_0,y_0)$，则$y_0^2=4x_0$。
			\par\textbf{理由：}切点未知，需引入参数$x_0,y_0$，且切点在抛物线上。
		\item[\textbf{第二步（写切线公式）：}] 由$y^2=4x$得$p=2$，故切线方程为$y_0 y = 2(x+x_0)$。
			\par\textbf{理由：}已知切点时可直接套用公式。
		\item[\textbf{第三步（代入已知点）：}] 切线过$A(4,5)$，代入得$5y_0 = 2(4+x_0)=8+2x_0$。
			\par\textbf{理由：}利用切线上点的坐标满足方程。
		\item[\textbf{第四步（联立解切点）：}] 由$y_0^2=4x_0$得$x_0=\dfrac{y_0^2}{4}$，代入$5y_0=8+2x_0$，得$5y_0=8+\dfrac{y_0^2}{2}$，整理得$y_0^2-10y_0+16=0$，解得$y_0=2$或$y_0=8$，对应$x_0=1$或$x_0=16$。
			\par\textbf{理由：}联立方程求解切点坐标。
		\item[\textbf{第五步（求切线方程）：}]
			当$(x_0,y_0)=(1,2)$时，切线$2y=2(x+1)$，即$y=x+1$；
			当$(x_0,y_0)=(16,8)$时，切线$8y=2(x+16)$，即$4y=x+16$。
			\par\textbf{理由：}每组切点对应一条切线，代回公式即可。
	\end{description}

	\textbf{\textcolor{CExample}{关键结论：}}
	所求切线有两条，分别为$y=x+1$和$4y=x+16$。

	\medskip
	\textbf{\textcolor{CExample}{例 3（综合应用）}}

	\textbf{\textcolor{CExample}{题目：}}
	已知抛物线$y^2=12x$的某条切线与直线$l: y=3x+2$平行，求该切线方程。

	\textbf{\textcolor{CExample}{解题步骤：}}
	\begin{description}[style=nextline, leftmargin=2.8em, labelsep=0.5em, itemsep=0.45em]
		\item[\textbf{第一步（求切线斜率）：}] 直线$l$的斜率为$3$，切线与之平行，故切线斜率$k=3$。
			\par\textbf{理由：}平行直线的斜率相等。
		\item[\textbf{第二步（由斜率求切点）：}] 对于$y^2=12x$，$2p=12$，$p=6$。由切线斜率公式$k=\dfrac{p}{y_0}$（来源于切线方程$y_0y=p(x+x_0)$变形），得$3=\dfrac{6}{y_0}$，解得$y_0=2$。代入抛物线得$4=12x_0$，解得$x_0=\dfrac{1}{3}$。故切点$\left(\dfrac{1}{3},2\right)$。
			\par\textbf{理由：}切线方程化为点斜式可得$k=\dfrac{p}{y_0}$，从而反求切点坐标。
		\item[\textbf{第三步（套用切线公式）：}] 将$y_0=2$，$x_0=\dfrac{1}{3}$，$p=6$代入$y_0 y = p(x+x_0)$得$2y = 6\left(x+\dfrac{1}{3}\right)$，即$2y = 6x+2$，化简得$y=3x+1$。
			\par\textbf{理由：}直接代入切点坐标得切线方程。
	\end{description}

	\textbf{\textcolor{CExample}{关键结论：}}
	所求切线方程为$y=3x+1$。
\end{examplebox}
